% --- 1. INFORMAZIONI GENERALI ---
\section{Informazioni Generali}
\begin{itemize}
	\item \textbf{Luogo/Piattaforma:} Server Discord
	\item \textbf{Data:} 2026-03-24
	\item \textbf{Orario di inizio:} 14:30
	\item \textbf{Orario di fine:} 15:22
	\item \textbf{Responsabile:} Sofia De Blasi
	\item \textbf{Scriba:} Francesco Marcon
\end{itemize}

\vspace{0.5cm}
\textbf{Partecipanti presenti:}
\begin{table}[ht]
	\centering
	\renewcommand{\arraystretch}{1.2}
	\begin{tabularx}{\textwidth}{X l}
		\toprule
		\textbf{Nome e Cognome} & \textbf{Matricola} \\
		\midrule
		Sofia De Blasi & 2111014 \\
		Roberto Mariano Doroftei & 2111031 \\
		Filippo Idiometri & 2035617 \\
		Francesco Marcon & 2110990 \\
		Armando Moda Scarati & 2082864 \\
		Anton Ricardo Rupi & 2054304 \\
		\bottomrule
	\end{tabularx}
\end{table}

\vspace{0.5cm}
\textbf{Partecipanti assenti:}
\begin{table}[ht]
	\centering
	\renewcommand{\arraystretch}{1.2}
	\begin{tabularx}{\textwidth}{X l}
		\toprule
		\textbf{Nome e Cognome} & \textbf{Matricola} \\
		\midrule
		Nessuno & - \\
		\bottomrule
	\end{tabularx}
\end{table}

% --- 2. ORDINE DEL GIORNO ---
\section{Ordine del Giorno}
Gli argomenti previsti per la riunione odierna sono:
\begin{enumerate}
	\item Revisione file per la candidatura.
	\item Aggiornamento sito web.
	\item Assegnazione prossimi verbali.
	\item Regole per repository (branch e Pull Request).
\end{enumerate}
\newpage
% --- 3. RESOCONTO ---
\section{Resoconto}

\subsection{Revisione file candidatura}
Il gruppo ha revisionato congiuntamente lo stato dei documenti per la candidatura, decidendo le seguenti azioni:
\begin{itemize}
	\item \textbf{Verbale esterno VarGroup:} Essendoci omonimia di date con altri file, andrà rinominato aggiungendo il suffisso \texttt{\_azienda}.
	\item \textbf{Norme di Progetto:} Sono state approvate formalmente.
	\item \textbf{Dichiarazione degli impegni:} È necessario uniformare la firma e il registro delle modifiche per renderli coerenti con gli altri documenti.
	\item \textbf{Valutazione dei capitolati:} Il file deve essere spostato nella cartella \texttt{doc\_esterni} e rinominato \texttt{valutazione\_capitolati}. Il registro delle modifiche va sistemato (versioni 0.0.0, 0.1.0 per le tabelle e 0.1.1 per correzioni ortografiche). Si è inoltre ricordato di verificare che i punteggi siano coerenti con la scala di preferenza stabilita e riflettano correttamente l'ordine finale di scelta dei capitolati.
	\item \textbf{Lettera di presentazione:} Andrà completata inserendo i link definitivi.
	\item \textbf{Glossario:} È stato approvato formalmente.
\end{itemize}
È stata inoltre discussa la definizione del parametro “complessità”, concordando di interpretarlo come adeguatezza rispetto alle capacità del gruppo.

\subsection{Aggiornamento sito web}
Si è stabilito che il sito web del gruppo dovrà essere aggiornato da Sofia, caricando i link definitivi a tutti i documenti richiesti per la consegna.

\subsection{Assegnazione prossimi verbali}
È stata definita l'organizzazione per i prossimi verbali: il verbale dell'attuale riunione interna viene assegnato ad Armando (con Roberto come revisore); il verbale dell'incontro di domani con l'azienda Zucchetti per i chiarimenti finali sul capitolato e la firma dei verbali esterni sarà redatto da Sofia (con Francesco come revisore).

\subsection{Regole per repository e documenti}
Riguardo al flusso di lavoro su GitHub, è stato imposto l'obbligo di creare un \textit{branch} dedicato per svolgere ogni task. Prima dell'apertura di una \textit{Pull Request}, si raccomanda di aggiornare il proprio branch rispetto al \textit{main} così da individuare e risolvere eventuali conflitti. Successivamente, la \textit{Pull Request} dovrà essere aperta assegnando esplicitamente la persona indicata come revisore.

\newpage
% --- 4. DECISIONI ---
\section{Decisioni}

\begin{table}[ht]
	\centering
	\renewcommand{\arraystretch}{1.5}
	\begin{tabularx}{\textwidth}{c X}
		\toprule
		\textbf{Codice} & \textbf{Decisione} \\
		\midrule
		\textbf{VI-5.1} & Approvazione formale dei documenti \textit{Norme di Progetto} e \textit{Glossario}. \\
		\textbf{VI-5.2} & Definizione delle correzioni necessarie sui documenti di candidatura (Valutazione Capitolati, Dichiarazione Impegni, Lettera di Presentazione e Verbali). \\
		\textbf{VI-5.3} & Assegnazione ruoli per la stesura e verifica dei prossimi verbali. \\
		\textbf{VI-5.4} & Adozione della regola obbligatoria sull'uso dei branch dedicati e delle Pull Request con revisore. \\
		\bottomrule
	\end{tabularx}
\end{table}

% --- 5. AZIONI DA SVOLGERE (TO-DO) ---
\section{Azioni da Svolgere (To-Do)}

\begin{table}[ht]
	\centering
	\renewcommand{\arraystretch}{1.5}
	\begin{tabularx}{\textwidth}{c c X l l}
		\toprule
		\textbf{Codice} & \textbf{Rif.} & \textbf{Descrizione Task} & \textbf{Assegnatario} & \textbf{Scadenza} \\
		\midrule
		\textbf{T-5.1} & VI-5.2 & Rinomina verbale VarGroup aggiungendo \texttt{\_azienda}. & Ricardo & 2026-03-24 \\
		\textbf{T-5.2} & VI-5.2 & Sistemazione firma e registro modifiche in Dichiarazione Impegni. & Armando & 2026-03-24 \\
		\textbf{T-5.3} & VI-5.2 & Spostamento in \texttt{doc\_esterni} e fix registro/punteggi di Valutazione Capitolati. & Filippo & 2026-03-25 \\
		\textbf{T-5.4} & VI-5.2 & Inserimento link definitivi nella Lettera di Presentazione. & Roberto & 2026-03-26 \\
		\textbf{T-5.5} & VI-5.2 & Aggiornamento del sito web con i link definitivi ai documenti approvati. & Sofia & 2026-03-26 \\
		\textbf{T-5.6} & VI-5.3 & Stesura verbale riunione esterna con Zucchetti. & Sofia & 2026-03-26 \\
		\textbf{T-5.7} & VI-5.4 & Creazione branch, lavoro e apertura PR con assegnazione revisore. & Tutti & 2026-03-26 \\
		\bottomrule
	\end{tabularx}
\end{table}

% --- 6. PROSSIMA RIUNIONE ---
\section{Prossima Riunione}
\begin{itemize}
	\item \textbf{Data:} {Da definire, presumibilmente primo giorno disponibile dopo l'aggiudicazione del capitolato.}
	\item \textbf{Ora:} {Da definire.}
	\item \textbf{OdG preliminare:} Analisi dell'esito dell'aggiudicazione dei capitolati; pianificazione delle attività di progetto in caso di esito positivo, oppure definizione delle strategie e dei prossimi passi operativi in caso di esito negativo.
\end{itemize}